\scanpage{5}
\deleted{Descente} \mspage{3}

% slanted
\margin{\textcircled{1} \uline{descente de propriétés de finitude}}

\uline{Prop} \uline{Soit} $S'\overset{f}{\mathrel{-}}S$ \uline{surjectif}. Pour qu'un $S$-morphisme $g:X\mathrel{-}Y$ soit surjectif (resp. radiciel) il f et s que $g':X'\mathrel{-}Y'$ le soit. [Si $f'$ inj $\mathrel{-}$ bij $\mathrel{-}$

\uline{Prop} \deleted{Soit} $f:S'\mathrel{-}S$ surjectif qu. cpt \inserted{(resp. fid plat qu. cpt)}. Alors $f$ est quasi-compact \inserted{ou quasi-séparé, i.e. $\Delta_f$} (resp. de t.f. \unsure{)}, loc. de t.f., de prés finie, loc de prés. finie) sss $f'$ l'est

\diagram{\begin{tikzcd}[ampersand replacement=\&,column sep=small]
\& X' \arrow[d] \& \& X \arrow[d,dash]\\
\unsure{Y'} \arrow[r,dash]\arrow[dr,dash] \& Y' \arrow[dr,dash] \& \unsure{X'} \arrow[r,dash]\arrow[dl,dash] \& X \arrow[r,dash]\arrow[d,dash] \& Y\arrow[dl,dash]\\
\& S'\arrow[rr,dash] \& \& S
\end{tikzcd}}

La partie non \unsure{resp.|respée} immédiate \deleted{\illegible{1}} On est ramené pour la partie \unsure{resp.|respée} : \unsure{au}.

\uline{Lemme} Soit $A\mathrel{-}A'$ fid plate, $B$ $A$-algèbre. Alors $B$ est de t.f. \inserted{(p.f.)} sss $B'$ l'est.

Les \unsure{$\Rightarrow$|mêmes} arguments, par modules, \unsure{donc|donnent}

\uline{Cor.} Soit $f:S'\mathrel{-}S$ fid plat qu. cpt, $\unsure{\mathcal F|F}$ faisceau qu. ch. sur $S$. Alors $\mathcal F$ de t.f. (pr. f.) sss $\mathcal F'$ l'est.

On prouve $\unsure{=|\text{même}}$ chose pour \unsure{libre|loc.}. libre de t.f. \inserted{loc libre de rang $n$ \quad donc : inversible}

\margin{\uline{Prop} $f$ fid plat. Alors $g$ est plat (fid plat) sss $g'$ l'est \quad ou \unsure{par modules|on peut prendre modules}}

% slanted
\margin{\textcircled{2} \uline{descente de propriétés topologiques} et conséquences\ldots}

Th.

\uline{Cor} \inserted{Soit $S'\mathrel{-}S$ fid plat \inserted{(ou univ. gén.)} et quasi-compact} Soit $f:X\longrightarrow Y$ un \unsure{morphisme|S-morphisme}. Alors \deleted{\unsure{Alors}} \inserted{\unsure{Si|$S'$}} $f'$ est ouvert (resp \inserted{générisant}, fermé, resp. quasi-compact et un homéo dans, resp. un homéo), $f$ de m.

Donc $f$ est imm ouvert (resp. \unsure{--}) sss $f'$ l'est.

\uline{Cor} $f$ séparé (propre) sss $f'$ l'est

\diagram{\begin{tikzcd}[ampersand replacement=\&]
 X'\arrow[r]\arrow[d,"g'"'] \& X\quad U\arrow[d,"g\ \text{surj}"]\\
 Y'\arrow[r,"f_g\ \text{surj}"'] \& Y
\end{tikzcd}}

\diagram{\begin{tikzcd}[ampersand replacement=\&]S'\arrow[r,dash,"{\text{surj}}","f"'] \& S\end{tikzcd}}

% slanted
\margin{\uline{Quelques} \uline{autres propriétés}}

\uline{Cor} $f$ un iso (mono) sss $f'$ l'est.

On est ramené au cas d'un iso, \deleted{et pour cela on est ramené} \inserted{On sait déjà par ce qui \unsure{pr.}} \deleted{\unsure{au fait} que si $f'$ l'est, $f$ est \unsure{a priori} ($f$ l'étant), donc} est un homéomorphisme, reste à prouver que $O_Y\mathrel{-}f_*(O_X)$ est un iso, or cela le devient après chgt. de base, o.k.

\deleted{\unsure{Cor} Soit $f$ fid plat, \illegible{1}}

\deleted{\uline{Question} Si $g'$ une imm, $g$ une ? O.K. si $g$ \unsure{loc. de prés. finie} ?}

surj, radiciel, \deleted{\unsure{quasi-sép}} \deleted{\unsure{injectif}} \inserted{plat, fid plat} qu cpt, qu. sép, \inserted{\unsure{cohér|coh.}}, t.f., p.f. loc t.f., loc p.f. séparé, propre, \deleted{\illegible{1}} univ. fermé, univ. ouvert, homéo univ, \unsure{domin.} biration\unsure{nel}

\unsure{Diagonal} \unsure{imm|immersions} (imm. fermée, imm. ouverte)

\deleted{\illegible{1}} affines, entiers, finis, quasi-affines, quasi-finis \deleted{plat fid}

[lisse, net, étale formellement lisse (?), formellemt étale\ldots.

\uline{TSVP}

% archive sheet number: 4
